\subsection{Isotropy groups of a cohomotoper}

Let G and H be groupoids and let \(\varphi\) , \(\psi\) be groupoid morphisms from H to G. The purpose of this section is to compute the isotropy groups of the cohomotoper \(\mathrm{Coh}(\varphi, \psi)\). We retain the notation \(G_1, h_1, \theta_1, \alpha, h\) of No. 3.

Denote by \(\Gamma_{0}\) the frame of the pair \((\varphi,\psi)\). Recall (II, p. 185, def. 3) that this is the quiver \((\mathrm{Orb}(\mathrm{G}),\mathrm{Orb}(\mathrm{H}),\varphi_{0},\psi_{0})\) , where \(\varphi_{0}\) and \(\psi_{0}\) are the maps from \(\mathrm{Orb}(\mathrm{H})\) to \(\mathrm{Orb}(\mathrm{G})\) induced by the maps \(\varphi\) and \(\psi\) by passage to orbits.

Throughout the rest of this section, we will further assume that the quiver \(\Gamma_{0}\) is connected and nonempty; by II, p. 185, prop. 4, this amounts to assuming that the groupoid \(\mathrm{Coh}(\varphi,\psi)\) is transitive, or equivalently that the quiver \(G_{1}\) is connected and nonempty (II, p. 185, prop. 2).

DEFINITION 4. --- We call base equipment of the pair \((\varphi, \psi)\) the data:

\begin{enumerate}
\def\labelenumi{(\roman{enumi})}
\item
  For every \(i \in \mathrm{Orb}(\mathbf{G})\), a vertex \(a(i)\) in the orbit \(i\) of \(\mathbf{G}\) ;
\item
  For every \(j \in \mathrm{Orb}(\mathrm{H})\), a vertex \(b(j)\) in the orbit \(j\) of H;
\item
  For every \(j \in \text{Orb(H)}\), arrows \(c_{1}(j)\) and \(c_{2}(j)\) in G connecting respectively \(\varphi(b(j))\) to \(a(\varphi_{0}(j))\) and \(\psi(b(j))\) to \(a(\psi_{0}(j))\) ;
\item
  A subquiver T of the frame \(\Gamma_{0}\) whose associated graph is a maximal tree of the graph \(\widetilde{\Gamma}_{0}\) ;
\item
  An orbit \(i_{0}\) of G.
\end{enumerate}

Let us choose a base equipment \((a,b,c_{1},c_{2},\mathrm{T},i_{0})\) of the pair \((\varphi,\psi)\). We define a quiver morphism \(\tau_{1}\) from \(\Gamma_{0}\) to \(\operatorname{Grp}(\mathrm{G}_{1})\) by setting \(\tau_{1}(i)=a(i)\) for \(i\in\operatorname{Som}(\Gamma_{0})=\operatorname{Orb}(\mathrm{G})\) and

\[
\tau_ {1} (j) = c _ {1} (j) ^ {- 1} \cdot h _ {1} (b (j)) \cdot c _ {2} (j)
\]

for \(j \in \mathrm{Fl}(\Gamma_{0}) = \mathrm{Orb}(\mathrm{H})\). We shall denote by \(\tau_{0}\) the composite of \(\tau_{1}\) and the canonical morphism \(\theta_{1}\) from \(\mathrm{Grp}(\mathrm{G}_{1})\) to \(\mathrm{Coh}(\varphi, \psi)\) ; it is a quiver morphism from \(\Gamma_{0}\) to \(\mathrm{Coh}(\varphi, \psi)\) .

For \(i \in \mathrm{Orb}(\mathrm{G})\) , denote by \(\alpha_i \colon \mathrm{G}_{a(i)} \to \mathrm{Coh}(\varphi, \psi)_{a(i)}\) the group homomorphism deduced from the morphism \(\alpha \colon \mathrm{G} \to \mathrm{Coh}(\varphi, \psi)\) by restriction to the isotropy groups at \(a(i)\) .

For \(j\in \mathrm{Orb(H)}\) , denote by

\[
\varphi_ {j} = \mathrm{Int} (c _ {1} (j)) ^ {- 1} \circ \varphi_ {b (j)} \colon \mathrm{H} _ {b (j)} \to \mathrm{G} _ {a (\varphi_ {0} (j))}
\]

and

\[
\psi_ {j} = \mathrm{Int} (c _ {2} (j)) ^ {- 1} \circ \psi_ {b (j)} \colon \mathrm{H} _ {b (j)} \to \mathrm{G} _ {a (\psi_ {0} (j))},
\]

so that one has, for every element \(f\) of \(\mathrm{H}_{b(j)}\) ,

\[
\varphi_ {j} (f) = c _ {1} (j) ^ {- 1} \varphi (f) c _ {1} (j) \quad \text { and } \quad \psi_ {j} (f) = c _ {2} (j) ^ {- 1} \psi (f) c _ {2} (j).\tag{6}
\]

For every vertex i of \(\Gamma_{0}\) , denote again by \(d_{i}\) the unique class of paths connecting \(i_{0}\) to i in the tree \(\widetilde{T}\) ; it is considered as an arrow in \(\mathrm{Grp}(\Gamma_{0})\) . Let us then denote by \(\delta_{i}\) the arrow of \(\mathrm{Coh}(\varphi,\psi)\) that is the image of \(d_{i}\) under the canonical morphism from \(\mathrm{Grp}(\Gamma_{0})\) to \(\mathrm{Coh}(\varphi,\psi)\) deduced from \(\tau_{0}\) ; the origin of \(\delta_{i}\) is \(a(i_{0})\) , its target is \(a(i)\) .

The quiver morphism \(\tau_0\) , the group homomorphisms \(\alpha_i\) (for \(i \in \mathrm{Orb}(\mathbf{G})\) ), \(\varphi_j\) and \(\psi_j\) (for \(j \in \mathrm{Orb}(\mathbf{H})\) ), and the arrows \(\delta_i\) in \(\mathrm{Coh}(\varphi, \psi)\) (for \(i \in \mathrm{Orb}(\mathbf{G})\) ) will be said to be deduced from the base equipment.

If \((\mathrm{G}_{i})_{i\in\mathrm{I}}\) is a family of groups, we denote by \(*G_{i}\) their free product; the image of an element \(g\in G_{i}\) under the canonical map from \(G_{i}\) into \(*G_{i}\) will be denoted {[}g{]}, or even g if there is no possible confusion. If S is a set, we denote by F(S) the free group constructed on S (A, I, p. 84).

PROPOSITION 6. --- There exists a unique group homomorphism

\[
\Lambda \colon \big (\underset {i \in \operatorname{Orb} (\mathrm{G})} {*} \mathrm{G} _ {a (i)} \big) * \mathrm{F} (\operatorname{Orb} (\mathrm{H})) \to \operatorname{Coh} (\varphi , \psi) _ {a (i _ {0})}
\]

such that

\begin{enumerate}
\def\labelenumi{(\arabic{enumi})}
\setcounter{enumi}{6}
\tightlist
\item
\end{enumerate}

\[
\Lambda (f) = \delta_ {i} \alpha_ {i} (f) \delta_ {i} ^ {- 1} \quad \text {   for   } i \in \operatorname{Orb} (\mathrm{G}) \text {   and   } f \in \mathrm{G} _ {a (i)},\tag{8}
\]

\[
\Lambda (j) = \delta_ {\varphi_ {0} (j)} \tau_ {0} (j) \delta_ {\psi_ {0} (j)} ^ {- 1} \qquad \text {for } j \in \mathrm{Orb} (\mathrm{H}).
\]

The homomorphism \(\Lambda\) is surjective; its kernel is the smallest normal subgroup of \((\ast_{i}\mathrm{G}_{a(i)})\ast\mathrm{F}(\mathrm{Orb}(\mathrm{H}))\) containing the elements \(j\) of \(\mathrm{Fl}(\mathrm{T})\) and the elements \(\varphi_{j}(f)j\psi_{j}(f)^{-1}j^{-1}\) , for \(j\in\mathrm{Orb}(\mathrm{H})\) and \(f\in\mathrm{H}_{b(j)}\) .

The existence and uniqueness of the homomorphism \(\Lambda\) follows from the universal property of free products and free groups (A, I, p. 85, prop. 8).

Let A be the set of \(a(i)\) for \(i \in \mathrm{Orb}(\mathrm{G})\) and \(\mathrm{G}_{\mathrm{A}}\) the full subgroupoid of G whose set of vertices is A. For every \(x \in \text{Som(G)}\) , denote by \(\overline{x}\) the orbit of \(x\) in G and choose an arrow \(d_x\) of G connecting \(x\) to \(a(\overline{x})\) . The pair \(v\) consisting of the map \(x \mapsto a(\overline{x})\) from \(\text{Som(G)}\) into A and the map which associates to \(f \in \text{Fl}_{x,y}(\text{G})\) the element \(d_x^{-1}fd_y\) of \(\text{Fl}_{a(\overline{x}),a(\overline{y})}(\text{G}_\text{A})\) is a morphism of groupoids. It follows from prop. 1 of II, p. 182 that \(v\) is a homotopism. Let \(\varphi' = v \circ \varphi\) and \(\psi' = v \circ \psi\) , and let \(w\) denote the canonical morphism from \(\text{Coh}(\varphi, \psi)\) to \(\text{Coh}(\varphi', \psi')\) ; it is a homotopism (II, p. 187, theorem 1).

The orbits of \(G_{A}\) are the sets \(\{a\}\) , for \(a \in A\) , and the injection \(G_{A} \to G\) induces a bijection from \(\mathrm{Orb}(G_{\mathrm{A}})\) onto \(\mathrm{Orb}(G)\) by which we will identify these two sets. One defines a base equipment \((a', b', \beta_{1}', \beta_{2}', T', i_{0})\) of the pair \((\varphi', \psi')\) by setting \(a'(i) = a(i)\) for \(i \in \mathrm{Orb}(G)\) , \(b'(j) = b(j)\) , \(\beta_{1}'(j) = v(c_{1}(j))\) , \(\beta_{2}'(j) = v(c_{2}(j))\) for \(j \in \mathrm{Orb}(H)\) and \(T' = T\) . The group homomorphisms \(\varphi_{j}'\) and \(\psi_{j}'\) (for \(j \in \mathrm{Orb}(H)\) ), the quiver morphism \(\tau_{0}'\) , the arrows \(\delta_{i}'\) (for \(i \in \mathrm{Orb}(G)\) ), and hence the group homomorphism \(\Lambda'\) , deduced from this base equipment are the composites, with w, of the corresponding homomorphisms \(\varphi_{j}\) , \(\psi_{j}\) , of the quiver morphism \(\tau_{0}\) , of the corresponding arrows \(\delta_{i}\) and of the homomorphism \(\Lambda\) .

Let B be the set of \(b(j)\) for \(j \in \mathrm{Orb}(\mathrm{H})\) , and let \(\mathrm{H}_{\mathrm{B}}\) be the full subgroupoid of H with vertex set B; denote by \(u: \mathrm{H}_{\mathrm{B}} \to \mathrm{H}\) the canonical injection; set \(\varphi'' = \varphi' \circ u\) and \(\psi'' = \psi' \circ u\) . The morphism \(u\) induces a bijection \(\mathrm{B} \to \mathrm{Orb}(\mathrm{H})\) by which we will identify these two sets. We again deduce from Theorem 1 of II, p. 187 a canonical homotopy equivalence \(w': \operatorname{Coh}(\varphi'', \psi'') \to \operatorname{Coh}(\varphi', \psi')\). Moreover, the pair \((\varphi'', \psi'')\) is equipped with a base equipment \((a'', b'', \beta_1'', \beta_2'', T'', i_0)\) , such that \(\Lambda', \varphi_j', \psi_j', \tau_0', \delta_i' (i, i' \in \operatorname{Orb}(\mathrm{G}), j \in \operatorname{Orb}(\mathrm{H}))\) are the composites with \(w'\) of \(\Lambda'', \varphi_j'', \psi_j'', \tau_0'', \delta_i'\) .

We summarize by the following diagram the various morphisms of groupoids introduced:

\[
\begin{array}{c} \mathrm {H_ {B}} \xrightarrow [ \psi^ {\prime \prime} ]{\varphi^ {\prime \prime}} \mathrm {G_ {A}} \xrightarrow {\alpha^ {\prime \prime}} \mathrm{Coh} (\varphi^ {\prime \prime}, \psi^ {\prime \prime}) \\ \Big \downarrow \qquad \qquad \qquad \Big \| \qquad \qquad \qquad \Big \downarrow w ^ {\prime} \\ \mathrm{H} \xrightarrow [ \psi^ {\prime} ]{\varphi^ {\prime}} \mathrm {G_ {A}} \xrightarrow {\alpha^ {\prime}} \mathrm{Coh} (\varphi^ {\prime \prime}, \psi^ {\prime \prime}) \\ \Big \| \qquad \qquad \qquad \Big | _ {v} \qquad \qquad \qquad \Big | _ {w} \\ \mathrm{H} \xrightarrow [ \psi ]{\varphi} \mathrm{G} \xrightarrow {\alpha} \mathrm{Coh} (\varphi , \psi)  . \end{array}\tag{9}
\]

To prove the proposition, we may therefore assume that \(A = \text{Som}(G)\) and \(B = \text{Som}(H)\) , in other words that the canonical maps \(\text{Som}(G) \to \text{Orb}(G)\) and \(\text{Som}(H) \to \text{Orb}(H)\) are bijective, assumptions under which we will proceed in the remainder of the proof.

The quiver \(G_{1}\) then has A as its set of vertices and the disjoint union of the sets \(G_{a}\) , \(a \in A\) , and B as its set of arrows. The arrows of \(G_{a}\) are loops at a; if \(b \in B\) , the arrow b connects \(\varphi(b)\) to \(\psi(b)\) , and the arrows \(c_{1}(b)\) and \(c_{2}(b)\) are loops at \(\varphi(b)\) and \(\psi(b)\) , respectively. The quiver T will be identified with an oriented tree of \(G_{1}\) ; it is a maximal oriented tree because the set of its vertices is equal to the set of vertices of \(G_{1}\) (II, p. 157, prop. 1). We set \(a_{0} = a(i_{0})\) . The set of arrows of the frame \(\Gamma_{0}\) of the pair \((\varphi, \psi)\) being identified with B, the quiver morphism \(\tau_{1}: \Gamma_{0} \to \text{Grp}(G_{1})\) associates to the arrow b the path class \(c_{1}(b)^{-1}bc_{2}(b)\) in the graph \(\widetilde{G_{1}}\) .

Recall that \(\theta_{1}\) denotes the canonical quiver morphism from \(\mathrm{G}_{1}\) to \(\mathrm{Grp}(\mathrm{G}_{1})\) . Denote by

\[
\lambda \colon \underset {a \in \mathrm{A}} {*} \mathrm{F} (\mathrm{G} _ {a}) * \mathrm{F} (\mathrm{B}) \to \operatorname{Grp} (\mathrm{G} _ {1}) _ {a _ {0}}
\]

the unique group homomorphism such that one has

\[
\begin{array}{l l} \lambda (f) = \tau_ {1} (d _ {a}) \theta_ {1} (f) \tau_ {1} (d _ {a}) ^ {- 1} & \text {if } a\in A \text{ and } f\in G_{a} \text{;} \\ \lambda (b) = \tau_ {1} (d _ {\varphi (b)}) \tau_ {1} (b) \tau_ {1} (d _ {\psi (b)}) ^ {- 1} & \text {if } b\in B \text{.} \end{array}
\]

Thus we have \(\lambda = \lambda' \circ \varepsilon\) , where \(\lambda'\) denotes the canonical group homomorphism from \(*_{a \in A} F(G_a) * F(B)\) to \(\text{Grp}(G_1)_{a_0}\) defined by the maximal oriented tree T (II, p. 179, prop. 9) and where \(\varepsilon\) is the unique automorphism of the group \(*_{a \in A} F(G_a) * F(B)\) such that \(\varepsilon(f) = f\) for \(a \in A\) and \(f \in G_a\) , and \(\varepsilon(b) = c_1(b)^{-1} h_1(b)c_2(b)\) for \(b \in B\) . By Remark 2 of II, p. 179, the homomorphism \(\lambda\) is surjective, and its kernel is the smallest normal subgroup of \(*_{a \in A} F(G_a) * F(B)\) which contains the arrows of T.

Denote by \(\pi\colon\mathrm{Grp}(\mathrm{G}_{1})\to\mathrm{Coh}(\varphi,\psi)\) the canonical groupoid morphism. By II, p. 177, cor. 1 of prop. 8, the group homomorphism \(\pi_{a_{0}}\) from \(\mathrm{Grp}(\mathrm{G}_{1})_{a_{0}}\) to \(\mathrm{Coh}(\varphi,\psi)_{a_{0}}\) is surjective, and its kernel is the smallest normal subgroup of \(\mathrm{Grp}(\mathrm{G}_{1})_{a_{0}}\) which contains the loops \(\mathrm{Int}(\tau_{1}(d_{a}))( \alpha_{1}(f)\alpha_{1}(g)\alpha_{1}(fg)^{-1})\) , for \(a\in\mathrm{A}\) and \(f,g\in\mathrm{G}_{a}\) and the loops \(\mathrm{Int}(\tau_{1}(d_{\varphi(b)}))( \varphi(f)b\psi(f)^{-1}b^{-1})\) , for \(b\in\mathrm{B}\) and \(f\in\mathrm{H}_{b}\) .

If \(p\colon F(\bigcup G_{a}\cup B)\to \left( \begin{array}{c} *\\ a\in A \end{array}G_{a}\right)*F(B)\) denotes the canonical surjective homomorphism, we thus have \(\Lambda \circ p = \pi_{a_0}\circ \lambda\) . This formula implies that the homomorphism \(\Lambda\) is surjective; it remains to determine its kernel.

For \(a \in \mathrm{A}\) and \(f \in \mathrm{G}_a\) , denote by {[}f{]} the image of \(f \in \mathrm{F}(\mathrm{G}_a)\) in the group \(*_{a \in \mathrm{A}} \mathrm{F}(\mathrm{G}_a) * \mathrm{F}(\mathrm{B})\) . For \(a \in \mathrm{A}\) , \(f, g \in \mathrm{G}_a\) , we then have

\[
\operatorname{Int} \left(\tau_ {1} \left(d _ {a}\right)\right) \left(\alpha_ {1} (f) \alpha_ {1} (g) \alpha_ {1} (f g) ^ {- 1}\right) = \lambda ([ f ] [ g ] [ f g ] ^ {- 1}).
\]

Similarly, for \(b \in \mathrm{B}\) and \(f \in \mathrm{H}_b\) the definition of homomorphisms \(\varphi_b\) and \(\psi_b\) (formula (6) of II, p. 193) implies that one has

\[
\begin{array}{r l} & {\mathrm{Int} (\tau_ {1} (d _ {\varphi (b)})) (\varphi (f) h _ {1} (b) \psi (f) ^ {- 1} h _ {1} (b) ^ {- 1})} \\ & {\quad = \tau_ {1} (d _ {\varphi (b)}) \varphi (f) \big (c _ {1} (b) \tau_ {1} (b) c _ {2} (b) ^ {- 1} \big) \psi (f) ^ {- 1}} \\ & {\qquad \qquad \big (c _ {2} (b) \tau_ {1} (b) ^ {- 1} c _ {1} (b) ^ {- 1} \big) \tau_ {1} (d _ {\varphi (b)}) ^ {- 1}} \\ & {\quad = \tau_ {1} (d _ {\varphi (b)}) c _ {1} (b) \varphi_ {b} (f) \tau_ {1} (b) \psi_ {b} (f) ^ {- 1} \tau_ {1} (b) ^ {- 1} c _ {1} (b) ^ {- 1} \tau_ {1} (d _ {\varphi (b)}) ^ {- 1}} \\ & {\quad = \lambda (c _ {1} (b)) \lambda (\varphi_ {b} (f) [ b ] \psi_ {b} (f) ^ {- 1} [ b ] ^ {- 1}) \lambda (c _ {1} (b)) ^ {- 1}.} \end{array}
\]

Consequently, the kernel of the homomorphism \(\pi_{a_0} \circ \lambda\) is the smallest normal subgroup of \(*\mathrm{F}(\mathrm{G}_a) * \mathrm{F}(\mathrm{B})\) which contains the elements \([f][g][fg]^{-1}\) for \(a \in \mathrm{A}\) and \(f\) , \(g \in \mathrm{G}_a\) , the elements \(\varphi_b(f)[b]\psi_b(f)^{-1}[b]^{-1}\) for \(b \in \mathrm{B}\) and \(f \in \mathrm{H}_b\) , and the elements \([b]\) , for \(b \in \mathrm{Fl}(\mathrm{T})\) .

Finally, the kernel of the homomorphism \(\Lambda\) is the smallest normal subgroup of \((\underset{a\in\mathrm{A}}{*}\mathrm{G}_{a})*\mathrm{F}(\mathrm{B})\) which contains the images under \(p\) of the preceding elements, in other words, the elements \([b]\) , for \(b\in\mathrm{Fl}(\mathrm{T})\) , and the elements \(\varphi_{b}(f)[b]\psi_{b}(f)[b]^{-1}\) , for \(b\in\mathrm{B}\) and \(f\in\mathrm{H}_{b}\) . The proposition is thus proved.

